\begin{solucion}
  \item Para este caso relizaremos la siguiente sustitución:
  \begin{align*}
    \alpha &\to \alpha^2\\
    \beta &\to \beta^2\\
    \gamma &\to \gamma^2\\
  \end{align*}
  Por lo tanto, el polinomio $P(y)$ se transforma en un polinomio $Q(y)$ que satisface:
  \begin{align*}
    e_1(\alpha^2, \beta^2, \gamma^2) &= \alpha^2 + \beta^2 + \gamma^2 = p\\
    e_2(\alpha^2, \beta^2, \gamma^2) &= \alpha^2\beta^2 + \alpha^2\gamma^2 + \beta^2\gamma^2 = q\\
    e_3(\alpha^2, \beta^2, \gamma^2) &= \alpha^2\beta^2\gamma^2 = r\\
  \end{align*}
  Estos polinomios son simetricos, por lo que podemos expresarlo en términos de los polinomios simétricos elementales,
  Por lo que expresando $\alpha^2, \beta^2, \gamma^2$ en términos de los polinomios simétricos elementales:
  \begin{align*}
    e_1(\alpha,\beta,\gamma), e_2(\alpha,\beta\gamma),e_3(\alpha,\beta,\gamma)
  \end{align*}
  Dado que $\mdeg(\alpha^2 + \beta^2 + \gamma^2) = (2,0,0)$. Por lo que tomamos:

  \begin{align*}
    g_0 &= e_1(\alpha,\beta,\gamma)^2\\
    &= (\alpha + \beta + \gamma)^2\\
    &= (\alpha + \beta)^2 + 2(\alpha + \beta)\gamma + \gamma^2\\
    &= \alpha^2 + \beta^2 + 2\alpha\beta + 2\gamma\alpha + 2\gamma\beta + \gamma^2\\
    &= \alpha^2 + \beta^2 + \gamma^2 + 2(\alpha\beta + \alpha\gamma + \beta\gamma)\\
  \end{align*}
  Por lo que:
  \begin{align*}
    \alpha^2 + \beta^2 + \gamma^2 &= (\alpha + \beta + \gamma)^2 - 2(\alpha\beta + \alpha\gamma + \beta\gamma)\\
    &= e_1^2 - 2e_2\\
    &= (-2)^2 - 2(-3)\\
    &= 4 + 6\\
    &= 10 =p
  \end{align*}
  Ahora para $(\alpha^2\beta^2 + \alpha^2\gamma^2 + \beta^2\gamma^2)$,podemos descomponerlo en polinomios simétricos 
  elementales 
  \begin{align*}
    e_1(\alpha\beta, \alpha\gamma,\beta\gamma), e_2(\alpha\beta, \alpha\gamma,\beta\gamma),e_3(\alpha\beta, \alpha\gamma,\beta\gamma)
  \end{align*}
  Asi:
  \begin{align*}
    \alpha^2\beta^2 + \alpha^2\gamma^2 + \beta^2\gamma^2 &= (\alpha\beta + \alpha\gamma + \beta\gamma)^2 - 2(\alpha^2\beta\gamma + \alpha\beta^2\gamma + \alpha\beta\gamma^2)\\
    &= e_2(\alpha,\beta,\gamma)^2 - 2(\alpha\beta\gamma)(\alpha + \beta + \gamma)\\
    &= e_2(\alpha,\beta,\gamma)^2  - 2e_3(\alpha,\beta,\gamma)e_1(\alpha,\beta,\gamma)\\
    &= (-3)^2 - 2(-2)(-5)\\
    &= 9 -20\\
    &= -11=q
  \end{align*}
  y por último:
  \begin{align*}
    \alpha^2\beta^2\gamma^2 &= (\alpha\beta\gamma)^2\\
    &= (-5)^2\\
    &= 25 = r
  \end{align*}

\end{solucion}